\chapter{Resolving singularities by surgery}
\label{ch:surgery}

This chapter describes one local surgery step and the continuation data that
make it usable in a changing-carrier Ricci flow.  A surgery is a geometric
operation only after its incoming slice, its strong necks, and its scale
parameters have been fixed.  We first construct the metric replacement on
those necks, then identify the topology of the resulting event and compare
its incoming and outgoing components.  The terminal-limit results of the
preceding chapter supply the necks; the global schedule chapter arranges
successive events.

\section{Changing carriers and terminal data}
\label{sec:surgery-carriers}

Let $F$ be a three-dimensional changing-carrier flow.  Its time domain is a
subset of $\mathbb R$, its slice at $t$ is a smooth manifold $M_t$, and its
metric and connection are denoted $g(t)$ and $\nabla(t)$.  Away from the
discrete set $\Sigma$ of surgery times, ordinary transport identifies the
slices and
\[
        \partial_t g(t)=-2\operatorname{Ric}(g(t)).
\]
At $T\in\Sigma$ the event has a pre-time $T^-$, a regular terminal region,
finitely many disjoint necks, a retained pre-region, a retention map into the
outgoing slice, and a finite family of inserted caps.  The incoming and
outgoing carriers are not identified globally.  All maps below are on the
part of a carrier where their domains are stated.

The flow data include compact slices, the ordinary transports, and the
event identifications.  Admissibility imposes three further conditions:
each boundary neck is strong; sufficiently late before a nonempty event,
every point outside the interior retained region has canonical control;
and sufficiently late before a vanishing event, every point has canonical
control.  Here canonical control means the neck, cap, or positive closed
component alternative with the flow's fixed constants $\varepsilon,C$.
The relevant late interval is part of each event.  Pinching means the
absolute-time Hamilton--Ivey inequalities displayed below on every slice.
A terminal
input at $T>0$ has time domain $[0,T)$, a last surgery-free slab
$[a,T)$, and curvature unbounded arbitrarily close to $T$.  It is
admissible, pinched, canonical above the prescribed radius threshold,
and noncollapsed with its prescribed function $\kappa(t)$.
The last condition is imposed on surviving parabolic balls with
backward interval $[t-r^2,t]$, $0<r\leq\varepsilon$, and
$|\operatorname{Rm}|\leq r^{-2}$: outside the positive-component
alternative their final volumes are at least $\kappa(t)r^3$.
The input also has a controlled core $\mathcal C\subset M_a$
and the scale data
\[
  \rho=\delta(T)r(T),\qquad 0<\rho<r(T),\qquad
  0<\rho<r_0,
\]
with
\[
 \delta(T)\leq\delta_0,\quad
 h(T)\leq R_0^{-1/2},\quad
 h(T)\leq\rho\delta(T),\quad
 h(T)\leq\frac{\rho}{2C}.
\]
The controlled core is classified by the disjoint alternatives
$\mathcal C\ne\varnothing$ and $\mathcal C=\varnothing$.
The terminal regular limit, the
selected deep horns, and the bridge from the last slab to that limit are
separate hypotheses of the continuation theorem.

\begin{definition}[Regular history]
\label{def:surgery-regular-history}
For a generalized reference flow $G$ on an interval $J$, a regular history
of $G$ in $F$ is a family of forward embeddings and inverse maps on every
ordinary box.  They are smooth, mutually inverse on the image, pull back the
metric exactly, and commute with the ordinary slab transports.  At a surgery
time they land in the interior of the retained region before and after the
event, and the retention map carries the pre-image to the same worldline in
the outgoing slice.  The history also contains the time-domain inclusion and
the compatibility of all overlapping boxes.
\lean{PoincareMT.M33RegularHistoryData}
\source{MT2007}
\dcref{Lemma 14.11 and Proposition 14.12, pp. 349--350}
\dcref{Definition 15.8, pp. 361--362}
\cite[Lemma 14.11 and Proposition 14.12, pp. 349--350;
Definition 15.8, pp. 361--362]{MT2007}
\end{definition}

The point of this definition is that a regular limit is not silently placed
on a fixed manifold.  The forward maps identify only the retained worldlines;
discarded neck ends have no image in the terminal regular region.

\section{The standard cap}
\label{sec:standard-cap}

\begin{definition}[Standard initial metric and maximal cap flow]
\label{def:standard-initial-metric}
A standard initial metric $g_0^{\rm cap}$ is a complete rotationally invariant
smooth metric on $\mathbb R^3$ with nonnegative sectional curvature and
sectional curvature $1/4$ on some ball about the origin.  It is invariant
under $\operatorname{SO}(3)$.  Outside a compact closed metric ball of radius
$A_0>0$, a specified collar identifies it with
$2g_{S^2}+ds^2$, where $g_{S^2}$ has sectional curvature one.
A partial standard-cap flow has lifetime $\tau>0$, solves Ricci flow on
$[0,\tau)$, and has both this initial metric and its specified connection.
For every $0\leq b<\tau$ there is $K_b<\infty$ bounding
$|\operatorname{Rm}|$ at every point and every time in $[0,b]$.
Completeness and rotational symmetry are conclusions, not additional
assumptions on a partial flow.  Maximality means that no flow with larger
lifetime agrees with both its metric and connection on its entire domain.
\lean{PoincareMT.StandardInitialMetric}
\source{MT2007}
\dcref{Definition 12.1, pp. 293--294}
\cite[Definition 12.1, pp. 293--294]{MT2007}
\end{definition}

\begin{lemma}[Common-domain uniqueness]
\label{lem:cap-common-uniqueness}
\uses{def:standard-initial-metric}
Two partial standard-cap flows with the same initial metric and connection
agree at every common included time, without a separate completeness
or symmetry assumption.
\lean{PoincareMT.M34.partialStandardCapFlow_metric_unique}
\source{MT2007}
\dcref{Section 12.5, pp. 309--319}
\cite[Section 12.5, pp. 309--319]{MT2007}
\end{lemma}

\begin{proof}
Fix a compact common time interval $[0,b]$.  Subtract the metrics,
connections, and raised curvature tensors of the two flows.  On the
compact tip and a sequence of translated cylindrical annuli choose
cutoffs whose plateaus cover the carrier and whose supports meet only
neighboring annuli.  Pull each translated annulus back to a fixed
reference domain, and denote the three tensor differences there by
$H_j,A_j,S_j$.  Fixed linear coordinates on their finite-dimensional
fibers give Euclidean norms.  With compactly supported cutoff $\chi_j$,
define the actual coordinate energy
\[
 E_j(t)=\int \chi_j^2
       \bigl(|H_j|^2+|A_j|^2+|S_j|^2\bigr)\,dx.
\]
Index zero uses a compact-core chart; positive indices use the
translated end chart.  Uniform metric comparison and derivative
estimates, including bounds inherited from the initial metric at
time zero, control the coefficients on these fixed supports.
Integration by parts in the curvature-difference
equation, together with the metric and connection evolution equations,
bounds $E_j'$ by a constant times the integral of the
same density over its support.  Derivatives of the cutoff give only
support terms: Young's inequality absorbs the factors containing
one derivative of the curvature difference into the negative
diffusion term.  Each support is covered by at most three adjacent
plateaus.  The uniform coefficient and cutoff bounds consequently give
$0\leq E_j\leq M$ with one $M$ for all $j$, and
\[
 E_0',E_1'\leq C(E_0+E_1+E_2),\qquad
 E_{j+2}'\leq C(E_{j+1}+E_{j+2}+E_{j+3}).
\]
All initial energies vanish.  To justify the noncompact argument, set
$F_k=\sum_{j=0}^{k+2}2^{-j}E_j$.  The displayed inequalities imply
\[
 F_k'\leq8CF_k+2CM\,2^{-(k+3)},\qquad F_k(0)=0.
\]
Scalar Gronwall bounds $F_k(t)$ by a quantity tending to zero as
$k\to\infty$, uniformly for $0\leq t\leq b$.  For each fixed $j$,
$2^{-j}E_j\leq F_k$ once $k\geq j$; hence $E_j=0$.
The plateau cover and continuity give metric equality everywhere;
uniqueness of the Levi-Civita connection gives connection equality.
Only finite sums were differentiated.  Exhausting the common time
interval proves the assertion.
\end{proof}

\begin{theorem}[Existence of the standard cap flow]
\label{thm:standard-cap-existence}
\uses{thm:neck-cap-canonical,thm:singular-m30}
\uses{def:standard-initial-metric,lem:cap-common-uniqueness}
For every standard initial metric $g_0^{\rm cap}$ there is a selected maximal
flow $G$ with $G(0)=g_0^{\rm cap}$ and lifetime exactly one.  For
$0\leq t<1$ every slice is complete and has nonnegative sectional curvature;
for $0<t<1$ it has positive sectional curvature and is rotation invariant.
For every $0\leq t_0<1$ and $\epsilon>0$, there is a compact set $K$ such
that each $x\notin K$ has one spatial cylinder chart of length
$\epsilon^{-1}$, valid for all $t\in[0,t_0]$, in which $G(t)$ is
$\epsilon$-close to $2(1-t)g_{S^2}+ds^2$.  Closeness uses spatial
covariant jets through order $\lfloor\epsilon^{-1}\rfloor$.
There are $r_*>0$ and $\kappa_*>0$ such that, whenever
$0<r\leq r_*$, $r^2\leq t<1$, and $|\operatorname{Rm}|\leq r^{-2}$
on the fixed final-time ball $B_{G(t)}(x,r)$ throughout $(t-r^2,t]$,
\[
 \operatorname{vol}_{G(t)} B_{G(t)}(x,r)\geq\kappa_*r^3.
\]
\lean{PoincareMT.repairedStandardCapExistence}
\source{MT2007}
\dcref{Theorem 12.5, pp. 295--297}
\dcref{Proposition 12.13, pp. 305--306}
\dcref{Theorems 12.28--12.29 and Claim 12.30, pp. 323--325}
\cite[Theorem 12.5, pp. 295--297; Proposition 12.13, pp. 305--306;
Theorems 12.28--12.29 and Claim 12.30, pp. 323--325]{MT2007}
\end{theorem}

\begin{proof}
First, the class of initial metrics is nonempty.  In radial arclength choose
a smooth concave profile $f$, equal to $2\sin(s/2)$ near zero, with
$0\leq f'\leq1$, and eventually constant $\sqrt2$.
A cutoff parameter adjusts the integral of $f'$ to $\sqrt2$; continuity
and the intermediate value theorem give the required plateau height.
The warped metric $ds^2+f(s)^2g_{S^2}$ has sectional curvatures
$-f''/f$ and $(1-(f')^2)/f^2$.  Thus it is nonnegative, is round with
curvature $1/4$ near the tip, and is the required cylinder at infinity.
The round formula makes the origin smooth, and the unbounded radial
arclength proves completeness.

For any specified initial metric, double successively longer truncated
cylindrical ends and run the compact Ricci flows.  Local curvature and
derivative bounds on fixed balls, together with the noncollapsing and
pointed-compactness results of the analytic chapters, give a diagonal
limit, including the initial endpoint.  On every compact time interval
with curvature bound $K$, integrating the Ricci equation gives uniform
two-sided comparison with the complete initial metric.  This proves
completeness of the limit.  The tensor maximum principle preserves
nonnegative curvature; its strong form, applied to the positive tip,
gives positive sectional curvature at positive times.
Lemma~\ref{lem:cap-common-uniqueness}, applied to rotation pullbacks, proves rotational
invariance.  Restarting and identifying overlapping solutions gives a
maximal flow.  Translating farther down the initial cylinder and taking
limits identifies the end with the shrinking cylinder, uniformly on
each compact time interval.  Reduced-geometry noncollapsing applies to
the controlled parabolic balls stated in the theorem.

The cap transfer also preserves the quantitative geometric margins used
later.  If a normalized core point has scalar curvature one and the
common cap constant is $C\geq1$, the reference ball of radius
$(100C)^{-1}$ has volume at least
$C^{-1}(100C)^{-3}$.  Forward and inverse chart differentials bounded
by $2$ put its image inside the corresponding target ball and change
calibrated volume by at most $2^3$.  Here Bishop--Gromov is applied in
the complete limit cap, whose Ricci curvature is nonnegative, to propagate
its core-ball volume witness down to the reference radius.  For a target
ball normalized by $\sup_{B(y,r)}R=r^{-2}$, the scalar bounds
$(2C)^{-1}\leq R\leq2C$ imply the conservative bounds
$(2C)^{-1}\leq r\leq2C$.  Short source paths and the differential
bound put the reference ball's image inside that target ball.  Consequently
\[
 \operatorname{vol}B(y,r)\geq
 \frac{(100C)^{-3}}{8C(2C)^3}\,r^3.
\]
This proves the required bound on the scalar-normalized core balls.
The same chart
bounds, applied to the coefficient derivatives on a larger compact ball,
give the finite spatial-jet bounds; they are not consequences of pointwise
metric convergence alone.

We give the essential lifetime argument because a bound on each compact
subinterval alone gives no bound at a finite endpoint.  If $\tau>1$,
the cylindrical asymptotic scalar curvature $1/(1-t)$ contradicts a
uniform curvature bound on a compact time interval containing $1$.
Suppose instead that $\tau<1$.  First there are constants $A,C,Q_*>0$
such that each point of scalar curvature at least $Q_*$ has canonical
control and $|\partial_tR|\leq AR^2$.  To prove this, choose the
requested neighborhood accuracy before its model constant.  If no
threshold worked, choose bad points with scalar curvature tending to
infinity.  Noncollapsing and the fixed-sequence long-limit theorem
give an ancient limit normalized to scalar curvature one at its base.
The ancient alternative is a strong neck or cap at that exact base.
For a neck, convergence on the backward unit slab transfers the
fixed cylinder chart.  For a cap, choose a relatively compact open
neighborhood of its closure, transfer both its end and boundary
neck charts, and retain the original core under the convergence
embedding.  Strict scalar-ratio, derivative, diameter, and volume
margins survive the finite-jet and volume convergence.  Completeness
allows the comparison balls to stay in the transferred neighborhood.
This produces a canonical neighborhood at the originally chosen
bad point.  Ancient analytic estimates, transferred through the
same space-time chart, give the scalar derivative bound there too,
a contradiction.  This is persistence of an ordinary limit cap,
and uses no theorem about caps born at a surgery event.

Because $\tau<1$, the asymptotic cylinder gives one finite scalar
bound $K$ at spatial infinity, uniform in the initial choice of
time $t_0<\tau$.  Choose $B\geq\max(1,Q_*,9K)$ first, then
$t_0=\tau-d$ with $0<d\leq\min(\tau/2,1/(16AB))$.
At that time some compact $X$ has $R(x,t_0)\leq B$ outside $X$.
For each fixed such $x$, the guarded inequality $R'\leq AR^2$
prevents $R$ from reaching $2B$ before $\tau$: integration of
$(1/R)'\geq-A$ on a last interval from $B$ to $2B$ would require
at least $1/(2AB)$ time.  Nonnegative sectional curvature bounds
$|\operatorname{Rm}|$ by scalar curvature.  We have therefore
obtained one exterior full-curvature bound on all of $[t_0,\tau)$.

Choose a relatively compact open set $\Omega$ outside $X$.  The exterior
curvature bound and metric comparison give
$\operatorname{vol}_{G(t)}\Omega\geq v>0$ for $t_0\leq t<\tau$.
Nonnegative Ricci curvature makes distances nonincreasing.  Compactness
of $X$ therefore supplies $D$ such that
$\Omega\subset B_{G(t)}(p,D)$ for every $p\in X$ and $t\geq t_0$.
Finite maximal lifetime forces curvature blowup, with centers
$p_j\in X$, times $t_j\uparrow\tau$, and scales $Q_j\to\infty$.
Point-picking and long-limit compactness give an ancient noncollapsed
limit of $(M,Q_jG(t_j),p_j)$; its asymptotic volume ratio is zero.
For fixed $A$ and large $j$, Bishop--Gromov comparison gives
\[
 \frac{v}{D^3}\leq
 \frac{\operatorname{vol}_{G(t_j)}B(p_j,D)}{D^3}
 \leq\frac{\operatorname{vol}_{Q_jG(t_j)}B(p_j,A)}{A^3}.
\]
Pass first to the pointed limit and then let $A\to\infty$.
The right side tends to zero, a contradiction.  Hence $\tau=1$.
\end{proof}

\begin{theorem}[Uniqueness, lifetime, and canonical alternatives]
\label{thm:standard-cap-uniqueness}
\uses{thm:standard-cap-existence,lem:cap-common-uniqueness}
Let $E$ be the existence data for $g_0^{\rm cap}$.  Every maximal standard-cap
flow $G'$ with the same initial metric has lifetime one and agrees with $E$ on
the common interval.  More generally, every compatible partial flow agrees
with $E$ on its common domain; maximality is needed only to compare lifetimes.
There is a constant $c>0$ such that, for every $t\in[0,1)$ and every
$x\in\mathbb R^3$,
\[
             \frac{c}{1-t}\leq R_E(x,t).
\]
For each $0<\epsilon<1/2$ there is $C>0$, chosen before $x$ and $t$, such
that every $(x,t)$ is either in a quantitative cap, is the center of an
evolving $\epsilon$-neck on $[0,t]$ disjoint from the initial cap core, or is
the center of an evolving neck on normalized time $(-(1+\epsilon),0]$.
More precisely, the cap alternative supplies an open $U\cong\mathbb R^3$,
a compact embedded closed ball $K\subset U$ with $x\in\operatorname{int}K$,
and an $\epsilon$-neck end $U\setminus K$.  Put $Q_U=\sup_U R$.
For $y,z\in U$ one has $0<R(z)<CR(y)$,
$d_U(y,z)<CQ_U^{-1/2}$, and
$\operatorname{vol}U<CQ_U^{-3/2}$.
For each $y\in\operatorname{int}K$ there is $r>0$ with compact
$\overline{B(y,r)}\subset U$, $\sup_{B(y,r)}R=r^{-2}$, and
$\operatorname{vol}B(y,r)>C^{-1}r^3$.
On $U$, $|\nabla R|<CR^{3/2}$ and $|\partial_tR|<CR^2$.
The evolving-neck alternatives use one fixed spatial chart and the
normalization $Q=R(x,t)$: the first has time interval $[-tQ,0]$
and misses the closed initial radius-$(A_0+4)$ ball; the second has
interval $(-(1+\epsilon),0]$.
\lean{PoincareMT.repairedStandardCapUniqueness}
\source{MT2007}
\dcref{Theorem 12.5, pp. 295--297}
\dcref{uniqueness argument, pp. 307--323}
\dcref{Proposition 12.31 and Theorem 12.32, pp. 325--327}
\cite[Proposition 12.31 and Theorem 12.32, pp. 325--327]{MT2007}
\end{theorem}

\begin{proof}
Common-domain uniqueness is Lemma~\ref{lem:cap-common-uniqueness}.
A maximal flow shorter than
the selected one extends by that flow; a longer one is excluded by the
same cylindrical lifetime barrier.

We establish the scalar rate before using the final longer-time neck
alternative.  The unit-time high-curvature model theorem gives
$A,H>0$ with $|\nabla R|\leq AR^{3/2}$ and
$|\partial_tR|\leq AR^2$ wherever $R\geq H$.
At a fixed time write the rotational metric in radial arclength as
$ds^2+w(s)^2g_{S^2}$.  Nonnegative curvature gives $w''\leq0$;
completeness and positivity for every $s>0$ force $w'\geq0$.
The end asymptotic gives an exterior scalar floor
$c_t=1/(2(1-t))$.  Since
\[
 R=\frac{2(1-(w')^2)}{w^2}-\frac{4w''}{w},
\]
one has $w^2\leq2/c_t=4(1-t)$ everywhere.  Indeed, if this
bound failed at some radius, monotonicity of $w$ would give a
strictly negative uniform upper bound for $w''$ farther out, forcing
$w'$ eventually negative.  Consequently, for a fixed nonzero point
$x$ and nonzero Euclidean tangent vector $v\perp x$,
\[
 G(t)_x(v,v)\leq4(1-t)\frac{|v|^2}{|x|^2}\longrightarrow0.
\]
If $R(x,t)$ did not diverge, the Harnack estimate would give a
uniform final scalar bound at $x$.  Nonnegative Ricci curvature and
the Ricci equation would then give a positive lower bound for
$G(t)_x(v,v)$, a contradiction.  To include the tip, use the
guarded gradient estimate on $R^{-1/2}$, clipped at the level $H$.
A bounded final scalar at the tip would bound scalar on a fixed small
metric ball around it.  Metric contraction makes a ball at time
$1/2$ stay inside that ball at all later times.  It contains a
nonzero point, contradicting the blowup just proved.

The reciprocal argument now gives the uniform rate.  In particular,
$\partial_tR\leq AR^2$ where $R\geq H$.
Put $\tau_0=\min(1/2,1/(2AH))$.  The terminal component of the set
$\{t:R(x,t)>H\}$ starts before $1-\tau_0$: otherwise integrate
$(1/R)'\geq-A$ from its left boundary to $1$ to obtain
$1/H\leq A\tau_0\leq1/(2H)$.
Thus, for $t\geq1-\tau_0$, integration to the blowup endpoint gives
$R(x,t)\geq1/(A(1-t))$.
On $[0,1-\tau_0]$ the compact tip and positive cylindrical asymptotic
give a uniform lower bound $B>0$.  Choosing
$c=\min(A^{-1},B\tau_0)$ proves the asserted estimate for all $x,t$.

Finally fix $\epsilon\in(0,1/2)$.  The high-curvature canonical
theorem supplies a threshold $r_*^{-2}$ and cap constants, with
neck limits extending for normalized time $1+\epsilon$.
The scalar rate just proved confines points below this threshold
to $t\leq\theta<1$.  On that interval the cylindrical asymptotic
gives the time-zero neck alternative outside a compact set, with
initial image avoiding the fixed cap core.  The remaining compact
space-time set has a common quantitative cap constant.  Taking the
maximum of these constants gives the stated $C$.
In deriving the longer neck interval the prescribed backward-time
compactness margin is used; an arbitrary positive extension would
not suffice for the requested $1+\epsilon$.
\end{proof}

\section{Metric surgery on a strong neck}
\label{sec:metric-surgery}

Let $N$ be a static $\delta$-neck with center scalar scale $Q$ and axial
coordinate $s\in(-\delta^{-1},\delta^{-1})$.  Write $A_0$ for the radius at
which the standard cap is exactly cylindrical.  At absolute time $t$, put
$\nu=\max\{-K_{\min},0\}$, where $K_{\min}$ is the least sectional
curvature on the neck.

\begin{theorem}[Metric surgery]
\label{thm:metric-surgery}
\uses{def:standard-initial-metric}
For every standard initial metric there are constants
\[
 0<\delta_0<1/200,\quad R_0>0,\quad q>1,\quad C_0>1
\]
chosen before the ambient manifold and the neck, with $100q<C_0$ after the
profile parameter is enlarged.  If $N$ has $\delta\leq\delta_0$, center scalar
curvature $Q\geq R_0$, and at time $t\geq0$ on its carrier
\[
 R\geq-\frac{6}{1+4t},\qquad
 R\geq2\nu\bigl(\log\nu+\log(1+t)-3\bigr)\quad(\nu>0),
\]
then there is a surgery output $(\widehat M,\widehat g)$ and a radial map
$\Psi:N\longrightarrow\widehat M$ with the following properties.

\begin{enumerate}
\item The negative half-neck, including its central sphere, is retained with
  its original metric; the positive tail is replaced by one closed standard
  cap.  The cap chart is defined on the open radius-$A_0+5$ ball, and its
  closed radius-$A_0+4$ ball is the closure of the inserted cap.
\item The metric is Hamilton--Ivey pinched, and the inner cap has positive
  sectional curvature.  If the whole input neck has positive sectional
  curvature, the whole output does as well.
\item $\Psi$ is continuous on the neck carrier, collapses a terminal part
  of the positive tail to the cap tip, and satisfies
  $d_{\widehat g}(\Psi x,\Psi y)\leq d_{g,N}(x,y)$,
  where $d_{g,N}$ is intrinsic distance in the input neck.
  Its inverse on the retained collar
  is smooth, and the retained metric identity holds on the closed
  nonpositive half-neck.
\item For every $\eta>0$ there is a further cutoff
  $\delta_{\rm cmp}(\eta)>0$ such that $\delta\leq\delta_{\rm cmp}(\eta)$
  gives a tip-preserving smooth chart on
  $B_{g_0^{\rm cap}}(0,\eta^{-1})$ whose image contains
  $B_{\widehat g}(\mathrm{tip},Q^{-1/2}\eta^{-1})$.
  In this chart the sum of squared spatial covariant-jet errors of
  $Q\widehat g$ relative to $g_0^{\rm cap}$, through order
  $\lfloor\eta^{-1}\rfloor$, is uniformly less than $\eta^2$.
\end{enumerate}
\lean{PoincareMT.repairedMetricSurgery}
\source{MT2007}
\dcref{Claim 13.1, Theorem 13.2 and Lemma 13.4, pp. 331--334}
\dcref{Corollary 13.12, pp. 339--340}
\cite[Theorem 13.2 and Lemma 13.4, pp. 332--334;
Corollary 13.12, pp. 339--340]{MT2007}
\end{theorem}

\begin{proof}
Choose $q$ so large that
\[
 100(A_0+4)^2<q,\qquad
 \frac{q^2}{s^4}e^{-q/s}<\frac1{100}
       \quad(0<s\leq A_0+4),
\]
and then choose $C_0>100q$.  In normalized coordinates put
\[
 f(s)=\begin{cases}C_0\delta e^{-q/s}&s>0,\\0&s\leq0.\end{cases}
\]
All derivatives vanish at zero.  On the output ball let
$s=A_0+4-\operatorname{dist}_{g_0^{\rm cap}}(0,\cdot)$.
With $\chi(s)=1$ for $s\leq5/4$ and $\chi(s)=0$ for $s\geq7/4$,
form the normalized preconformal metric
\[
 J=\chi\,Q(\Psi^{-1})^*g+(1-\chi)(1-6\delta)g_0^{\rm cap}.
\]
The first summand is used only on its smooth collar domain, and its
weight vanishes outside that domain.  Choose a positive round-tip
radius $r\leq A_0$ and put $A_*=A_0+4$.  A second cutoff $\zeta$
equals one for $s\leq A_*-3r/4$ and zero for $s\geq A_*-r/2$.
The physical metric is
\[
 \widehat g=Q^{-1}
 \bigl(\zeta(s)e^{-2f(s)}+(1-\zeta(s))e^{-2f(A_*)}\bigr)J.
\]
This multiplier is constant near the tip and equals $e^{-2f}$ on
the neck transition.  Thus the radial-coordinate singularity does
not spoil smoothness at the origin.

The curvature estimate uses the formula, for a $J$-orthonormal pair
$e_1,e_2$,
\[
 K_{e^{-2f}J}(e_1\wedge e_2)=e^{2f}
 \bigl(K_J+\nabla^2f(e_1,e_1)+\nabla^2f(e_2,e_2)
 +(df(e_1))^2+(df(e_2))^2-|df|^2\bigr).
\]
Here $f'=qf/s^2$ and $f''=(q^2/s^4-2q/s^3)f$.
The quantitative order of choices matters.  The near-cylinder geometry
first supplies, for the normalized blend $J$, the bounds
$1/2\leq|\nabla s|_J^2\leq2$, Hessian and Laplacian errors at most
$K\delta$, scalar curvature at least $1/2$, and sectional curvature
at least $-K\delta$. Moreover, a $J$-orthonormal plane spanned by $v,w$
has curvature at least $1/8$ whenever
$ds(v)^2+ds(w)^2\leq1/2$. This is the nearly transverse plane case.
For $\phi(s)=C_0\delta e^{-q/s}$ the profile estimates are chosen so that,
with $B=K\delta$ and $\kappa=1/8$,
\[
 2B\phi'(s)+2\phi'(s)^2\leq\tfrac14\phi''(s),\qquad
 2B\phi'(s)+2\phi'(s)^2<\kappa,\qquad
 2\nu_J\phi(s)\leq\tfrac14\phi''(s).
\]
The first inequality absorbs both Hessian and gradient-square errors in
the conformal sectional-curvature formula when the plane has a substantial
height component; the second preserves the transverse positive gap, and
the third controls conformal amplification of the negative-curvature part
$\nu_J$ of the blend. Choose
$C_0>\max\{100q,128K e^q/q^2\}$; after shrinking $\delta_0$,
the profile also satisfies $K\delta<\phi''/4$ for $1\leq s\leq2$.
Choosing $q$ first, then $C_0$, and only then shrinking
$\delta_0$ makes these inequalities uniform on $0<s\leq2$ (the factor
$e^{-q/s}$ controls the finitely many cutoff derivatives).  The
conformal estimate gives, for $0<s<2$,
\[
 R_{\widehat g}\geq Q/2,\quad R_{\widehat g}\geq Q R_J,
 \quad\nu_{\widehat g}\leq Q\nu_J.
\]
For $s\geq1$ the same gain inequality gives positive sectional curvature
on every orthonormal plane; on the pure-cap part this is combined with
the round-tip curvature estimate and the constant multiplier.  Thus the
smaller radius-$(A_0+3)$ cap is positive as well.  On the overlap where
the second cutoff is changing, its derivative terms are bounded by the
same finite profile table and are included in the preceding choice of
$C_0$ and $\delta_0$.
For $s\leq0$ both metric and curvature are unchanged. On $0<s<1$,
the blend is exactly the normalized pullback of the old metric, so here
$R_{\widehat g}\geq R_g$ and $\nu_{\widehat g}\leq\nu_g$.
For $s\geq1$, positivity gives $\nu_{\widehat g}=0$ directly.
To deduce logarithmic pinching on $0<s<1$, set
$H_t(v)=2v(\log v+\log(1+t)-3)$.
If $H_t(\nu_{\widehat g})\leq0$, nonnegative scalar curvature suffices.
Otherwise $H_t$ is increasing between $\nu_{\widehat g}$ and $\nu_g$,
so $R_{\widehat g}\geq R_g\geq H_t(\nu_g)
\geq H_t(\nu_{\widehat g})$.  This avoids assuming that $H_t$ is
increasing on all positive numbers.

The radial collapse equals the original identification on the negative
half and becomes constant at the tip.  The quadratic-form estimates
for the contracted blend give $|d\Psi(v)|_{\widehat g}\leq|v|_g$
on each smooth piece.  Partitioning a piecewise smooth neck path at
the transition collars, and passing by approximation across the
constant locus, gives the length inequality.  Taking infima over neck
paths proves the asserted distance bound.  On the closed negative
half the same argument in both directions gives intrinsic isometry.
An ambient path leaving the neck is not among these paths.

Finally each fixed finite set of profile and cutoff derivatives is
bounded, while the normalized input neck tends to the cylinder as
$\delta\to0$.  The transition-coordinate estimates control the metric,
connection, and curvature jets on the closed radius-$(A_0+4)$ ball;
the pure-cap chart is constant near the tip, so no radial-coordinate
singularity enters this estimate.  The error in the cap metric therefore
tends to zero in each prescribed finite spatial jet norm.  Choosing the
further cutoff after $\eta$ gives the last conclusion.  The lower and
upper radial distance estimates imply respectively that the physical
image contains the ball of radius $Q^{-1/2}\eta^{-1}$ and that the chart
is defined on the standard radius-$\eta^{-1}$ ball.  Thus the comparison
is on the actual metric ball and not merely on a coordinate neighborhood.
\end{proof}

\section{Continuation through a singular time}
\label{sec:surgery-continuation}

The continuation input consists of a preterminal slab, its generalized reference flow,
the regular terminal limit and the selected horns, and a bridge carrying the
same metric and scalar data between the slab and the reference.  The bridge
also supplies the Appendix A neck--cap topology at accuracy
$c_T\varepsilon$, where $c_T=10^{14}$, with
$c_T\varepsilon\leq\varepsilon_A$ for the supplied neck--cap theory.
The bridge identifies the reference extension with the horn selection's
limit, preserves metric, scalar curvature and slab transport, and takes
the controlled core exactly onto the terminal sublevel
$\{R_T\leq\rho^{-2}\}$.  It identifies the constants
$\varepsilon,C,r_0$ and supplies metric surgery with the same chosen
constants on each selected neck.  Its deep-horn rule applies to every
calibrated strong horn at the stated accuracy and returns a neck at
the prescribed height $h(T)$, not at a newly selected height.

\begin{theorem}[Branch-correct continuation]
\label{thm:branch-continuation}
\uses{thm:singular-m31,cor:singular-selector}
\uses{def:surgery-regular-history,thm:metric-surgery}
Assume compact Ricci-flow local existence, uniqueness and continuation,
and the absolute-time pinching theorem.  Every regular-history window has
regular-history data.  Given a terminal input $I$ at $T$, a singular regular
limit $L$, horn data $N$, and a bridge satisfying the metric pullback,
scalar pullback, transport, core-image, calibration, and metric-insertion
hypotheses, there is a continuation extension on the same input.  It
preserves all earlier events, has the terminal cut-and-cap event when the
controlled core is nonempty, and has a permanent empty continuation when the
core is empty.  In either case the extension has a positive surgery-free
interval after $T$, remains admissible and pinched, and contains the supplied
regular history on its old domain.
\lean{PoincareMT.repairedBranchContinuation}
\source{MT2007}
\dcref{Lemma 15.11, p. 364}
\dcref{Theorem 15.9, pp. 363--366 and 409--410}
\cite[Lemma 15.11, p. 364; Theorem 15.9, pp. 363--366 and 409--410]{MT2007}
\end{theorem}

\begin{proof}
Ordinary slab transport first constructs compatible boxes along each
retained worldline.  Choose a countable refinement; on overlaps the
transport cocycle identifies the same points and tangent vectors.
The local metric jets therefore define a single smooth horizontal
metric on this atlas.  Across an old event only the continuing interior
is identified.  Equality of its limiting metrics and the Ricci equation
give equality of all one-sided time jets there.  This constructs the
regular history without identifying discarded worldlines.

Suppose first that the controlled core is nonempty.  Its image is a
compact scalar sublevel.  Since components of the regular limit are
open, a finite subcover proves that only finitely many components meet
it.  We must cut every escaping end of those same components.  When
$C\leq99$, the small-constant calibrated-horn construction gives the
required original-end horns directly.  For $C>99$, the global
neck--cap alternative gives either a calibrated horn or a capped tube.
If the tube contains a point with $R_T\leq\rho^{-2}$, cut off a
closed escaping half of that tube, keeping its original end.
Otherwise the low point lies in its cap.  The fixed-accuracy cap-overlap
lemma then excludes a canonical cap centered where
$R_T\geq2C\rho^{-2}$: overlap would force the low point's curvature
to be comparable to that center, contrary to this threshold.
Thus those high points are strong-neck centers.  Continuing overlapping
necks above $4C\rho^{-2}$ produces a horn whose boundary is below
$8C\rho^{-2}$ and whose tail is the original end.

Deep-horn selection now gives cuts at height $h(T)$ with controlled
boundary curvature, in particular the bound $32C\rho^{-2}$ used in
the overlap argument.  Whole neck collars, rather than just their
central spheres, lie in the selected ends.  Choose a maximal pairwise
disjoint family of candidate neck carriers.  All centers have scalar
curvature $h(T)^{-2}$, a compact level by scalar properness.  Uniform
neck geometry supplies a fixed positive packing radius there, so the
disjoint family is finite.  Maximality ensures that every candidate
intersects a selected neck.  This intersection alone does not give
end coverage: the calibrated boundary estimate and height inequalities
show that its original escaping end has a tail inside the selected
horn cut, even when the intersection is shallow.  For this finite
list, directed neck overlap makes intersecting discarded tails nested;
remove contained redundant tails.  Every original end remains covered.
An escaping sequence in the complement would determine an uncovered
end, contradicting that property.  Removing these
tails leaves a region with compact closure and disjoint spherical
boundary.  This compactness is needed for restart; finiteness of the
list alone would not imply it.

Apply Theorem~\ref{thm:metric-surgery} on these disjoint collars and
glue the cap balls to the retained region.  The collar identifications
are smooth and isometric on a neighborhood of the retained closed
half; they produce one compact smooth outgoing manifold and one
pinched metric.  Compact Ricci-flow existence restarts it at $T$.
Uniqueness joins overlapping ordinary restarts to a maximal one,
while the absolute-time pinching theorem retains the original clock.
At the event the horizontal metrics glue on the continuing interior;
the newly exposed caps and the discarded ends are not a common smooth
slice across $T$.  Pulling back the compact retained set to one late
pre-time constructs the actual event data and preserves all old events.

If the controlled core is empty, every disappearing worldline either
has scalar curvature tending to infinity or has terminal scalar
greater than $\rho^{-2}>r(T)^{-2}$.  Compactness gives a common late
interval on which all disappearing points have canonical control.
Thus deleting the entire slice satisfies admissibility.  Record that
vanishing event and set every later slice to the empty manifold; all
flow and pinching conditions there are vacuous.  This branch has
infinite lifetime and no later event.  Empty controlled core does not
assert that the regular-limit carrier itself was empty.
\end{proof}

The continuation retains its supplied metric surgeries and all old
identifications; no separate compatibility construction is required.

\section{The topology of the actual event}
\label{sec:surgery-event-topology}

Metric insertion and topological reconstruction have different hypotheses.
The next statement requires the neck--cap topology theorem and the
admissibility of the supplied flow; it does not assume an already
classified list of discarded factors.

\begin{theorem}[Local reconstruction]
\label{thm:surgery-local-topology}
\uses{thm:neck-cap-cover,thm:metric-surgery}
For any supplied neck--cap theory with threshold $\varepsilon_A>0$,
there is $0<\varepsilon_0\leq\varepsilon_A$, chosen before the flow,
such that every admissible three-dimensional surgery flow with
$c_T\varepsilon\leq\varepsilon_0$, with the terminal accuracy factor
$c_T=10^{14}$, has the following property at each
event $T$.  Its incoming slice is obtained by a finite smooth
connected-sum assembly from the actual outgoing components, finitely
many $S^2$-bundles over $S^1$, and finitely many closed positive
spaceforms.  The assembly includes collared balls and gluing maps.
In a nonempty event every actual inserted cap lies in its recorded
survivor component.  In a vanishing event there are no survivor pieces.
\lean{PoincareMT.rawLocalSurgeryTopology}
\source{MT2007}
\dcref{Proposition 15.3, pp. 357--358}
\dcref{Appendix A.21}
\cite[Proposition 15.3, pp. 357--358]{MT2007}
\end{theorem}

\begin{proof}
Transport the finite central neck spheres into one sufficiently late
incoming slice.  Cut there and cap both sides.  On the retained side
the actual event charts identify the capped result with the outgoing
slice, including its particular cap balls.  On the other side form
the capped discarded carrier using the same collars.  The collars
give compatible smooth charts; disjointness and closed attachment
graphs make the quotient Hausdorff.  Its compactness follows from
compact discarded closures and finitely many closed cap balls.

Choose the late slice uniformly for these finitely many closures.
Admissibility supplies canonical neighborhoods at all their points.
The neck--cap continuation theorem places each closure inside a
canonical tube, capped tube, or closed canonical model.  The containing
model is not identified with the discarded component.  Write $U$ for
the image of that component in the containing model $Q$.  Its frontier is
the finite disjoint union of the actual incident collar spheres.
The positive sides of these collars enter $U$; their negative sides
miss $\overline U$.

First suppose each frontier sphere bounds a smooth closed ball $B_i$ in
$Q$.  Choose $p\in U$.  Since $U$ is connected and misses $\partial B_i$,
it lies on one side of that sphere.  If $p\in B_i$, then
$\overline U\subset B_i$, giving spherical coordinates around the whole
closure.  Otherwise $U$ misses every $B_i$.  Two fillings with disjoint
frontiers are nested or disjoint.  Nesting is impossible here: the inner
frontier belongs to $\overline U$, whereas the interior of the outer ball
misses $\overline U$.  The balls are therefore disjoint.  Their complement
is connected, and $U$ is both open and relatively closed in it, since all
of $\partial U$ lies on their frontiers.  Thus
\[
                 U=Q\setminus\bigcup_i B_i.
\]
Adjust the ball collars to the given event collars, preserving their
central spheres and side choices.  Gluing those adjusted balls back in
then identifies the literal capped event component with $Q$.  In the
enclosing-ball alternative, the same construction in spherical
coordinates identifies it with a sphere.

For a positive spaceform, the required frontier fillings are constructed
by lifting each sphere to its finite spherical cover.  Smooth sphere
filling upstairs gives a finite family of balls on sides avoiding a fixed
point.  An innermost ball has disjoint nontrivial deck translates, so the
covering is injective on it and its descent is a ball with the prescribed
frontier downstairs.  This supplies the preceding argument for closed
spaceforms and punctured projective caps.  A tube or Euclidean cap has
spherical region coordinates.  A capped tube first absorbs its attached
tube into its original cap by the relative cap-absorption result of the
neck--cap chapter.  Its smallness threshold is fixed before the flow.

There are two additional closed models.  In a sphere bundle over the
circle, either all frontier spheres fill and the preceding construction
returns that bundle or a sphere, or one sphere is essential.  Lift the
essential sphere to the pullback cylinder over $\mathbb R$ and straighten
its collar.  Cut along a slightly negative collar slice, disjoint from
$\overline U$.  Connectedness places the entire closure in one covering
sheet of the cut complement, which has spherical coordinates.  Its
capped component is consequently spherical.
For a projective double, the corresponding cylinder cover has an infinite
dihedral deck action.  If all frontier spheres fill, the exact-complement
alternative reconstructs the whole double and hence its two projective
summands; the enclosing alternative is spherical.  Otherwise straighten
an essential lifted collar and again cut at a negative slice.  A connected
lifted region either projects injectively, giving spherical coordinates,
or is preserved by a reflection, giving punctured-projective coordinates.
The spaceform filling construction then applies to the actual region.

Components incident to no collar are whole old components and use the
closed canonical classification directly.  These constructions classify
and assemble each literal capped discarded component, including all its
incident spheres, before the incoming slice is reconstructed.

Now append this assembly to the unchanged outgoing components and
reverse the cuts using their recorded balls.  The finite incidence
graph has vertices the retained and capped discarded components and
edges the event collars.  A spanning forest joins distinct pieces;
remaining identifications give the residual bundle factors.  In
particular, a separating sphere gives a connected-sum separation,
whereas a nonseparating sphere contributes a sphere-bundle factor.
Reversing the finite list yields a diffeomorphism onto the original
incoming slice.  Tracking the retained vertices preserves each
original cap's survivor assignment.  With no cuts, the construction
reduces to whole-component classification.  If the outgoing slice
is empty, that classification contains no retained vertices.
\end{proof}

\section{Comparison across one surgery}
\label{sec:comparison}

Let $T$ be a nonempty surgery event.  Choose compact connected parent
and child components $X,Y$, with their smooth inclusions $\iota_X$
and $\iota_Y$ in the incoming and outgoing slices.  Write
$K_{\rm pre}$ for the retained pre-region and $\mathfrak r$ for
retention.  A comparison input records the selected local reconstruction,
the actual parent and child pullback metrics, separating surgery
spheres, and the full retained open set
\[
 \mathcal R=\{x\in X:\ \iota_X x\in\operatorname{int}K_{\rm pre},\quad
       \mathfrak r(\iota_X x)\in\iota_Y(Y)\},
\]
and the fact that $\mathcal R\ne\varnothing$.  The event must satisfy
\[
 \delta(T)<\delta_0,\qquad h(T)<R_0^{-1/2},\qquad
 c_T\,\varepsilon\leq\varepsilon_0.
\]

\begin{theorem}[Comparison map and metric bounds]
\label{thm:comparison-map}
\uses{thm:metric-surgery,thm:surgery-local-topology}
There is $\varepsilon_0>0$, chosen before the standard metric and flow, such
that every comparison input above has a continuous map
\[
                   \Phi:M_{T^-}^{\rm parent}\longrightarrow M_T^{\rm child}
\]
which agrees with the retention map on $\mathcal R$, is smooth there, has
open retained target image, and sends every point outside $\mathcal R$ into
one of the actual closed caps.  On $\mathcal R$ its metric is the pullback of
the terminal limit metric.  For every $\eta>0$ there is
$0<d\leq T-T^-$ such that,
for $T-d<t<T$, the same map satisfies
\[
 d_{g_T}(\Phi x,\Phi y)\leq (1+\eta)d_{g(t)}(x,y)
 \quad\text{for all }x,y.
\]
The target basepoint is the image under this map of the chosen parent
basepoint.
\lean{PoincareMT.repairedComparisonMap}
\source{MT2007}
\dcref{Proposition 15.12, p. 365}
\dcref{Claim 18.22, p. 433}
\dcref{Claim 13.1 and Theorem 13.2, pp. 331--334}
\cite[Proposition 15.12, p. 365; Claim 18.22, p. 433]{MT2007}
\end{theorem}

\begin{proof}
Let $K$ be the connected closed retained root for the chosen child.
Its boundary consists of the relevant separating neck spheres.
For each boundary sphere choose the connected complementary side
containing the outward collar.  That side misses $K$: otherwise a
path in $K$ from the inward collar would cross the sphere from the
wrong side.  Two such outward sides are disjoint; if they met,
connectedness and the disjoint frontiers in $K$ would force one to
cross the other's frontier.  The root and all outward sides cover
the parent, since their union is open and closed there.

On the root use retention.  On each outward collar use the radial
collapse, and beyond its constant locus send the whole remaining
branch to that cap's tip.  The definitions agree on collar overlaps
and are constant near each outer gluing boundary, so they define a
continuous map.  On the full open retained region the map is smooth
with the stated metric identity.  All remaining images lie in the
actual closed caps.

Crucially, all nonconstant local models are controlled on one compact
subset $K_*$ of the regular region.  Smooth convergence there gives,
for all sufficiently late $t$, the quadratic-form inequality
$g_T^{\rm reg}\leq(1+\eta)^2g(t)$ on $K_*$.
At each parent point the map has one of three local descriptions:
constant, retained isometry, or neck collapse.  The first has zero
image length, the second uses this metric inequality, and the third
combines it with the local intrinsic-distance estimate of metric
surgery.  Partition any piecewise smooth parent path into finitely
many segments in these neighborhoods.  Sum the local inequalities
and take the infimum over paths to obtain the global distance bound.
The single compact control set gives one time window for every path
and pair of points.  Shrinking it to the available pre-event interval
gives the stated bound on $d$.
\end{proof}

\begin{theorem}[Degree, homotopy, and third homotopy transport]
\label{thm:comparison-transport}
\uses{thm:comparison-map,prop:foundations-top-homology,prop:foundations-cw-equivalence}
Assume in addition that the parent and child components are closed and simply
connected and that a parent integral orientation
\[
 H_3(M_{T^-}^{\rm parent};\mathbb Z)\simeq\mathbb Z
\]
has been fixed.  For the literal comparison map $\Phi$ just constructed there
are a child orientation and a degree-one identity on third homology.  The map
is the forward map of a homotopy equivalence and induces a bijection on the based
third homotopy groups.  For every $\eta>0$ and all sufficiently late
pre-times there is a smooth $(1+\eta)$-Lipschitz map $\Phi_t$ with the exact
target basepoint, homotopic to $\Phi$, having the same degree and the same
based $\pi_3$ map.
\lean{PoincareMT.repairedComparisonHomotopy}
\source{MT2007}
\dcref{Proposition 15.12, p. 365}
\dcref{cap-image discussion, p. 428}
\dcref{Poincare transport, pp. 430--431}
\dcref{Claim 18.22, p. 433}
\cite[Proposition 15.12, p. 365; Claim 18.22, p. 433]{MT2007}
\end{theorem}

\begin{proof}
Choose a point $x$ in the retained open set whose image $y$ lies outside
the closed caps.  Injectivity of retention and the cap-image condition
give $\Phi^{-1}(y)=\{x\}$.  Near $x$ the map is a local homeomorphism.
It consequently gives an isomorphism on the local groups
$H_3(X,X\setminus\{x\};\mathbb Z)$ and
$H_3(Y,Y\setminus\{y\};\mathbb Z)$, where $X,Y$ denote the parent
and child.  The restriction maps from global third homology to these
local groups are isomorphisms for connected closed oriented
three-manifolds, as proved in the foundations chapter
(the standard local-homology formulation is
\cite[Theorem 3.26(a), pp. 236--238]{Hatcher}).
The naturality square therefore makes
$\Phi_*:H_3(X;\mathbb Z)\to H_3(Y;\mathbb Z)$ an isomorphism.
Given the chosen isomorphism $o_X:H_3(X;\mathbb Z)\to\mathbb Z$,
set $o_Y=o_X\Phi_*^{-1}$.  Then $o_Y\Phi_*=o_X$, exactly the
degree-one identity, independently of an initial choice of child
orientation.

The foundations chapter supplies homotopy equivalences of each closed
simply connected three-manifold with $S^3$, constructed from its finite
CW model and homology.  Conjugating $\Phi$ by these equivalences gives
a self-map of $S^3$ inducing an isomorphism on $H_3$, hence of degree
$1$ or $-1$.  Its inverse homotopy class gives a homotopy inverse;
conjugating back makes the actual $\Phi$ a homotopy equivalence.
This uses the homotopy type of a sphere, not a prior homeomorphism
with a sphere.  The based third homotopy map is therefore bijective.

For completeness, near-isometric smoothing must reserve room both
for smoothing and for the basepoint correction.  On compact smooth
manifolds a finite chart cover permits local convolution with cutoff
interpolation.  Choose nested source charts and positive target-chart
margins.  At stage $j$, convolution changes the map by at most $e_j$
and its Lipschitz constant by at most $s_j$; choose these budgets with
summable totals below the desired displacement and Lipschitz errors.
Previously smoothed neighborhoods remain smooth.  Short geodesics in
the target charts give the homotopy at each stage.  A finite partition
of any path upgrades the local derivative bounds to a global
Lipschitz estimate.

Put $\alpha=\min(1,\eta/8)$.  First choose a late time window on
which $\Phi$ is $(1+\alpha)$-Lipschitz.  A supported target
diffeomorphism, homotopic to the identity and with Lipschitz constant
at most $1+\alpha$, can move any sufficiently close point to the
prescribed target basepoint.  Choose the smoothing displacement within
that radius and its Lipschitz increase at most $\alpha$.  Composing
the smooth approximation with this correction gives an exactly based
smooth map with Lipschitz constant
\[
 (1+\alpha)(1+2\alpha)\leq1+\eta.
\]
Its free homotopy to $\Phi$ preserves degree.  The possible trace
of the basepoint is a loop in the simply connected target, so its
action on $\pi_3$ is trivial.  Thus the based homomorphism is the
same, not merely conjugate under an unspecified basepoint change.
\end{proof}

\section{The two continuation branches}
\label{sec:branch-dichotomy}

\begin{corollary}[Terminal branch alternatives]
\label{prop:unified-branch}
\lean{PoincareMT.repairedBranchContinuation}
\uses{thm:branch-continuation}
For the terminal input and continuation constructed in
Theorem~\ref{thm:branch-continuation}, a nonempty controlled core gives
a nonempty outgoing slice and the prescribed neck surgeries. An empty
core gives an empty terminal slice and an extension by empty slices
for all later times, with no later surgery. Both alternatives retain
the same old history and event identifications.
\source{MT2007}\dcref{Lemma 15.11, p. 364}
\end{corollary}
\begin{proof}
Inspect the controlled-core alternative in the constructed continuation.
In the nonempty case the terminal operation cuts the selected necks and
attaches the standard caps. In the empty case the construction continues
by empty carriers. Admissibility and pinching on these carriers are
vacuous; the old-event preservation is unchanged.
\end{proof}

\section{Persistence of a surgery cap}
\label{sec:cap-persistence}

An observation has a horizon $H>0$, interval $[0,H)$, and a selected
standard cap flow.  Fixing scales after a time $a$ means that the
initial cap model, metric-surgery constants, $\varepsilon,C$, and
height selector are those of the fixed parameter setup, and that,
for $a\leq s<H$,
\[
 r_{\rm next}\leq r(s),\quad \delta(s)\leq\bar\delta,
 \quad h(s)=h_{\rm sel}(\delta(s)r(s),\delta(s)).
\]
Canonical control at radius $r_{\rm next}$ means that every point
with $R\geq r_{\rm next}^{-2}$ has the specified strong canonical
neighborhood.  These conditions permit $r(s)$ to change at an epoch
boundary; they do not assert that the height is constant in time.

\begin{theorem}[Cap persistence or removal]
\label{thm:cap-persistence}
\uses{thm:metric-surgery,thm:standard-cap-uniqueness}
Fix a standard initial metric $g_0^{\rm cap}$, its selected standard
flow $G$, and the metric-surgery
constants.  Given a finite initial list of epoch parameters, ending with
radius $r_i$, and a next radius $r_{\rm next}$
with
\[
 0<r_{\rm next}\leq r_i,
\]
and positive $A,\eta,\theta$ with $\theta<1$, there is a cutoff
$\bar\delta>0$ chosen before the actual flow, observation, event, and cap.
For every admissible pinched flow whose observation uses that same
standard flow, assume the fixed-scale conditions after the preceding
epoch's start $a$ and canonical control at radius $r_{\rm next}$
on $[0,H)$.  Let $a\leq T<H$ be a nonempty surgery event with
$\delta(T)\leq\bar\delta$.  For each inserted cap with tip $p$, put
\[
 h=h(T),\quad U=B_{g(T)}(p,Ah),\quad
 T_{\rm end}=\min(H,T+\theta h^2),\quad
 d=(T_{\rm end}-T)/h^2.
\]
Then one of the following holds.

\begin{enumerate}
\item The worldlines of all of $U$ give a flow cylinder for normalized
  time $s\in[0,d)$.  A tip-preserving initial chart maps
  $B_{g_0^{\rm cap}}(0,A)$ exactly onto $U$; the cylinder is the
  identity on $U$ at time zero.  Pull back by this fixed spatial
  chart after normalizing time by $h^2$ and metric by $h^{-2}$.
  Relative to $G(s)$, the sum of squared spatial covariant-jet
  errors through order $\lfloor\eta^{-1}\rfloor$ is bounded by
  one number strictly less than $\eta^2$, uniformly in $x,s$.
\item There is a later surgery $T<T_+<T_{\rm end}$ and the same
  comparison on $[0,(T_+-T)/h^2)$.  The whole tracked ball disappears
  there: either the outgoing slice is empty, or, sufficiently late
  before $T_+$, every tracked worldline lies outside the interior
  retained region of that event.
\end{enumerate}

The initial chart is linked to the metric-surgery cap chart by the actual
local insertion map.  The normalization is performed once, before
pulling back both tensor slots.  The cutoff is chosen before the later
flow, observation, event, and cap.
\lean{PoincareMT.repairedCapPersistence}
\source{MT2007}
\dcref{Proposition 16.5, pp. 370--371}
\dcref{Claim 16.6, Corollary 16.7,
Lemma 16.8, Corollary 16.9, pp. 371--373}
\dcref{Claim 16.10, pp. 374--375}
\cite[Proposition 16.5, pp. 370--371; Claim 16.6 and Corollary 16.9,
pp. 371--373; Claim 16.10, pp. 374--375]{MT2007}
\end{theorem}

\begin{proof}
Suppose the conclusion fails for every cutoff and choose counterexamples
with $\delta_j\to0$.  Normalize each birth cap by its own $h_j^{-2}$
and translate its birth time to zero.  The incoming strong neck exists
for one backward normalized time unit.  Interior derivative estimates
therefore bound every curvature derivative on its middle half at
birth.  The fixed surgery cutoffs and smooth cap metric transfer these
bounds to the newly inserted region.  In particular, this is an
initial-endpoint estimate; positive-time smoothing alone would not
justify convergence at the birth of a cap.

Choose a buffered collar and a larger fixed standard ball
$B_{A+b}$ containing the closed radius-$A$ ball, with $b>0$ fixed
before the sequence.  The initial strong-neck derivative estimates and
the cap transfer give uniform bounds for the metric, connection, and all
spatial jets needed on $B_{A+b}$ at birth.  Canonical scalar derivative
bounds, together with pinching and local curvature transport, then give
the corresponding bounds for a definite normalized time.  In coordinates
the Ricci equation is an ODE for the metric coefficients and the
connection equation is an ODE for their first derivatives; integrating
the common bounds gives a uniform modulus in time.  Choose the time
increment so that this modulus is smaller than the collar margin $b$.
The chart consequently continues on the same worldlines and the image
of the physical closed ball remains inside the buffered chart.  These
are local estimates on the tracked buffer, not a scalar-rate bound on
every point of the whole surgery flow.

The exact birth-ball chart requires a further construction. Write
$\widehat g_j=h_j^{-2}g_j(T_j)$ and fix a large parameter radius $R>2(A+b)$.
The initial comparison confines closed $\widehat g_j$-balls in a larger
compact chart. Its metric and connection coefficients converge smoothly
to those of $g_0^{\rm cap}$. Choose orthonormal frames at the actual tips.
After a subsequence their coordinate matrices converge to an orthonormal
model frame $L$. The geodesic equation and its differentiated equations
give convergence, on fixed parameter balls, of the actual exponentials
$E_j$ to the standard exponential $E_0$ in that frame.

The standard tip is a pole: radial arclength gives a global smooth
exponential and inverse logarithm. In its logarithmic coordinates the maps
$E_j$ have derivatives converging to the identity on each smaller closed
parameter ball. A derivative error less than one, integrated along a line
segment in this convex ball, proves injectivity; nonsingularity gives a
smooth inverse. These are actual exponential charts. Each radial geodesic
has initial norm $|v|$, hence length $|v|$ to parameter one. Conversely,
compact containment of a physical ball gives a minimizing geodesic from the
tip to every point in it. Initial-value uniqueness identifies it with the
radial curve of $E_j$, with parameter norm equal to its distance. Therefore
\[
 E_j(B(0,r))=B_{\widehat g_j}(p_j,r)
 \qquad(0<r\leq A+b).
\]
The map $C_j=E_j\circ E_0^{-1}$ consequently sends every standard
radius-$r$ ball exactly onto this physical normalized ball, for the same
map $C_j$. In the original comparison coordinates its adjustment tends
smoothly to the identity. The bilinear pullback formula, with the derivative
in both tensor slots and one factor $h_j^{-2}$, preserves convergence of
every fixed finite jet. Thus the adjusted charts still satisfy the required
strict comparison error. A hypothetical failure at every positive cutoff
would give such a convergent sequence, contradicting the resulting charts;
the adjustment is therefore uniform before the physical flow is selected.

Compose with the local insertion embedding. Its metric pullback is exact,
and compact ball containment prevents shortcuts through the complement of
its image, so the images are the ambient birth balls in the surgery slice.
Scaling distances gives
$B_{\widehat g_j}(p_j,r)=B_{g_j(T_j)}(p_j,h_jr)$.
Apply the time-modulus and compactness estimates on the larger closed ball
$B_{A+b}$ first, and only then restrict to $B_A$, retaining these exact
birth identities. A diagonal sequence of larger balls and pointed
compactness give a partial standard-cap flow.  Completeness follows
from uniform metric comparison on each compact time interval.
Theorem~\ref{thm:standard-cap-uniqueness} identifies this limit with
$G$.  Since $\theta<1$, the curvature of $G$ is bounded on the
relevant compact time interval.  The local estimates can therefore
be restarted by a positive amount depending on this bound and the
fixed control constants.  Repetition gives convergence throughout
the assigned interval wherever the tracked cylinder survives.
The same adjusted initial chart is used throughout the comparison cylinder.

It remains to show that a loss of one worldline removes the whole
tracked ball.  If its maximal cylinder stops before $T_{\rm end}$,
bounded curvature and ordinary local continuation force that endpoint
to be a later surgery $T_+$.  Enlarge the test ball first if necessary
so that it contains a standard neck sphere bounding a ball about
the tip.  Suppose the tracked ball met a later cutting sphere.
If $h'$ is that surgery's scale, its scalar curvature is comparable
to $(h')^{-2}$, whereas comparison with $G$ bounds curvature on the
tracked ball by $Kh^{-2}$.  Thus $h/h'$ is bounded.  The tracked
ball's diameter is bounded by a multiple of $Ah$; the later strong
neck has buffer length of order $(\delta')^{-1}h'$.  For sufficiently
small cutoff, the entire tracked ball therefore lies inside that
neck.

The neck-overlap sphere lemma makes its internal standard neck
sphere homotopic to the central sphere of the later neck.  The
latter is nontrivial in $\pi_2(S^2\times I)$, while the former
bounds its tip ball inside the tracked region, hence inside that
neck.  This is impossible.  Therefore the tracked ball avoids all
later cutting spheres near $T_+$.  It is connected, so once one of
its worldlines lies on a disappearing side all of them do.
The comparison already proved holds up to $T_+$, giving the second
alternative.  Both alternatives contradict the chosen counterexamples;
one positive cutoff consequently works uniformly as asserted.
\end{proof}

\section{From one event to a controlled history}
\label{sec:surgery-interfaces}

The constructions above preserve one set of event identifications.
The terminal-limit theorem supplies the regular region and calibrated
horns; metric surgery supplies its cap metric; local reconstruction
records the actual retained and discarded factors.  The comparison
map uses those same factors and metrics.  Its based $\pi_3$ transport
can therefore be used for loop-family classes without changing the
underlying surgery event.

To iterate this operation one must still select scale parameters,
propagate canonical neighborhoods and noncollapsing, and rule out
finite accumulation of event times.  Cap persistence supplies the
local input for that induction.  The global controlled-flow chapter
performs the parameter selection and proves the finite-event bounds;
the later finite-history argument reverses the exact assemblies
constructed here.
